% The two-regime estimate that pays one regular source atom.
\section{A paid estimate for one regular source}
\label{sec:ex-lattice-trace-regular-atom}

We spell out how a single regular source term is paid in the selected
one-block transition. The source is split into its real and imaginary
channels. When the transported support fits inside the receiving budget
radius, local coefficient probes control the inserted coefficient. When it
does not fit, a direct supremum estimate controls it. The two cases use
different fields of the local budget, which is why both the regular and
large-support inequalities were selected in
Lemma~\ref{lem:ex-lattice-local-budget-choice}. The principal insertion
attenuation needed in the local estimate is established in
Section~\ref{sec:regular-principal-producer}; the argument below keeps the
regular source and both support regimes explicit.

Fix a one-step transition of block length \(L\) between the adjacent
rungs \(m\) and \(m+1\) of the selected history, at fixed volume and period
indices, and a coefficient frame of
maximum grade \(d\). Both ends of the transition are the effective
representations of that one history at these indices; the source, its
mask, the cloud estimate and the resulting coefficient all refer to this
same pair. Let \(f\) be a complex function of the fine (rung \(m+1\))
configuration and \(b\) a fine support such that:
\begin{enumerate}
\item \(f\) is regular on \(b\) for the fine budget, that is, local on
\(b\), gauge invariant, measurable, and smooth to the order of the budget
on configurations healthy on \(b\);
\item the origin belongs to \(b\), and \(f(\mathrm{id})=0\);
\item if \(b\) is nonwrapping about the origin, then for every local jet
the function \(s\mapsto f\) along the corresponding logarithmic-link test
curve at the origin is \(C^d\) at \(s=0\) and all its derivatives of order
at most \(d\) vanish there;
\item there is \(c\geq0\) such that, on every configuration healthy on
\(b\), every ordered derivative of \(f\) of order at most the order of
the fine budget, along directions satisfying the fine budget's
directional bound, has absolute value at most \(c\).
\end{enumerate}
Condition (4) with the empty list of directions is the pointwise bound
\(|f|\le c\) on healthy configurations; the full derivative bound, and
conditions (1)--(3), are needed for the local estimates below, which do
not follow from a value bound. The source term itself is the product of
the fine effective density with the mask of \(f\) on \(b\) at the fine
threshold. Let \(\alpha_f>0\) be the fine support
exponent and \(q_f\geq0\) the fine cover number of \(b\), and define the
charged mass
\[
 m_0=e^{\alpha_f q_f}c.
\]
Let \(\widehat\beta_m>0\) be the hatted coupling of the history at the
coarse rung \(m\) (not the actual coupling of the finite-volume
representation), let
\(\mathcal M_f(G,L,\widehat\beta_m)>0\) be the fine monitor, and put
\(\chi=\sqrt{\mathcal M_f(G,L,\widehat\beta_m)}\). Thus a raw budget \(b_0\)
gives the actual channel price \(b_0/\chi\). Set \(a=L^d\). Let
\(O,H,W\geq0\) be respectively the motion-average norm, local readout,
and local test factors. Let raw budgets \(u,v,\xi>0\) satisfy
\begin{align}
 O(Wu+2Hav)&\leq \xi, \label{eq:ex-lattice-regular-small-budget}\\
 4OHav&\leq \xi. \label{eq:ex-lattice-regular-large-budget}
\end{align}
These are the regular and large-support paid budgets. In the Wilson
application, \(m_0\) is the mass of the source support with its sup-bound;
\(u,v\) are chosen by Section~\ref{sec:ex-lattice-local-paid-budgets}.

The local estimates needed from the one-step construction have the following
precise form. The real and imaginary channels have transported cloud
charges \(C_1,C_2\) satisfying
\begin{equation}
 0\leq C_1,C_2\leq \frac{v}{\chi}m_0. \label{eq:ex-lattice-regular-cloud}
\end{equation}
In the source both charges are the single cloud response of the
two-channel source, so \(C_1=C_2\), and \eqref{eq:ex-lattice-regular-cloud}
is the cloud entry of the selected local record at price \(v\); it is a
separate input, not a consequence of the principal reserved-mass estimate
of Section~\ref{sec:regular-principal-producer}.
If the transported support is at most the receiving budget radius, the
local probe reconstruction error, after both channels are combined, is at
most
\begin{equation}
O\bigl(Wu\chi^{-1}m_0+Ha(C_1+C_2)\bigr). \label{eq:ex-lattice-regular-small-read}
\end{equation}
If the transported support exceeds that radius, each channel admits a
uniform chart bound \(S_i\leq \frac{v}{\chi}m_0\) for \(i=1,2\); the resulting two-channel
coefficient error is at most
\begin{equation}
 OHa(S_1+S_2+C_1+C_2). \label{eq:ex-lattice-regular-large-read}
\end{equation}
Here the \emph{coefficient error} is measured uniformly over the finite
set of lattice motions, and \(O\) pays for this uniformity. In either
case the local construction also supplies the following certificate, in
which \(\rho=L^{-2/5}\) is the one-step remainder rate. There is one
coarse coefficient \(X\), chosen before any motion, such that for every
motion \(\omega\) there is a coarse presentation \(P_\omega\) with
coefficient \(\omega X\) which has the order-zero format of the coarse
representation, has remainder cost at most \(\rho m_0\) at the
normalization \(\sqrt{\operatorname{forcedAllowance}_m}\), and has the
exact pushed density
\begin{equation}
 D(P_\omega)=\omega\cdot\Bigl[L^d\,\mathcal T_m\bigl(
     \varrho_{m+1}\,\operatorname{mask}_b f\bigr)\Bigr],
 \label{eq:ex-lattice-regular-pushed-density}
\end{equation}
where \(\varrho_{m+1}\) is the fine effective density,
\(\operatorname{mask}_bf\) the mask of \(f\) on \(b\) at the fine
threshold, \(\mathcal T_m\) the one-step push at the coupling of rung
\(m+1\), and \(\omega\cdot\) the action of the motion on functions. The
presentation may depend on the motion; the coefficient \(X\) does not.
The estimates \eqref{eq:ex-lattice-regular-small-read} and
\eqref{eq:ex-lattice-regular-large-read} are bounds for
\(\sup_\omega\|\omega X\|\), the predicted leading coefficient of a regular
source being zero. These
are the actual local chart and transport conclusions used below, not
consequences of the scalar budget inequalities alone.

\begin{proposition}[Regular source atom]
\label{prop:ex-lattice-regular-source-atom}
Assume the local estimates \eqref{eq:ex-lattice-regular-cloud}--
\eqref{eq:ex-lattice-regular-large-read} and the coarse-format,
density, and remainder conclusions just stated. If
\eqref{eq:ex-lattice-regular-small-budget}--
\eqref{eq:ex-lattice-regular-large-budget} hold, then the source term
produces one coarse coefficient \(X\) with
\(\|\omega X\|\le\xi\chi^{-1}m_0\) for every motion \(\omega\), together
with presentations \(P_\omega\) of \(\omega X\) in the coarse order-zero
format having the
exact pushed density \eqref{eq:ex-lattice-regular-pushed-density} and
remainder at most \(\rho m_0\). In the terminology
of the finite price recurrence, this is the regular-source estimate with
cost pair \((\xi\chi^{-1}m_0,\rho m_0)\).
\end{proposition}
\begin{proof}
First suppose the transported support is within the receiving radius. By
\eqref{eq:ex-lattice-regular-cloud},
\[
O\bigl(Wu\chi^{-1}m_0+Ha(C_1+C_2)\bigr)
\leq O(Wu+2Hav)\chi^{-1}m_0\leq \xi\chi^{-1}m_0,
\]
where the last inequality is \eqref{eq:ex-lattice-regular-small-budget}.
The local probe reconstruction estimate
\eqref{eq:ex-lattice-regular-small-read} therefore gives the claimed
coefficient error.

Now suppose the transported support exceeds the receiving radius. The
uniform chart estimate and the two channel bounds give
\[
 OHa(S_1+S_2+C_1+C_2)
 \leq OHa(4v\chi^{-1}m_0)=4OHav\chi^{-1}m_0\leq \xi\chi^{-1}m_0,
\]
by \eqref{eq:ex-lattice-regular-large-budget}. Thus the large-support
estimate gives the same coefficient error. The two cases exhaust the
support sizes. In both cases, the local construction supplies the exact
pushed-density identity, coarse format, and remainder bound \(\rho m_0\);
these components are unchanged by the numerical comparison. This proves the
assertion.
\end{proof}

Here is the source realization of the local hypotheses. The source term is
turned into its real and imaginary masked channels. Both channels carry the
same root label, the logarithmic key of the regular activity with support
\(b\). Their two reserved-mass bounds are combined at the single expansion
index \(i_0\) attached to that label, where the assigned cost is
\(L^d(c_{\mathrm{re}}+c_{\mathrm{im}})\) with \(c_{\mathrm{re}},
c_{\mathrm{im}}\) the derivative bounds of the two transported channels;
every other expansion index is
assigned cost zero. The per-channel reserved-mass estimate used for this
format and remainder certificate is
provided by the regular insertion result of
Section~\ref{sec:regular-principal-producer}; it does not supply the
cloud bound \eqref{eq:ex-lattice-regular-cloud} or the large-support chart
bound below, which use two further entries of the selected local record. In the small-support regime,
the finite local probes fit inside the chart, and the local reconstruction
bound is the test-price term \(Wu\chi^{-1}m_0\) plus the two readout terms
\(HaC_1\) and \(HaC_2\). In the large-support regime, the direct chart
supremum bound replaces the probes, leaving two chart and two cloud
readout contributions; this is exactly the factor four in
\eqref{eq:ex-lattice-regular-large-budget}. The cloud inequality controls
each \(C_i\), and the insertion construction carries the same source
through the block with its density identity. The finite price recurrence
then sums these single-source costs over the countable regular family using
its absolute mass budget.

More explicitly, let \(q_\tau\) be the coefficient in
Section~\ref{sec:ex-lattice-local-paid-budgets}. In the small-support case,
the chart estimate gives, for each recovery probe \(\tau\), a contribution
at most \(q_\tau v\chi^{-1}m_0\). The cycle budget gives
\(q_\tau v\leq u\); finite coefficient recovery therefore bounds the
probe part by \(Wu\chi^{-1}m_0\). The two real/imaginary cloud
contributions are bounded by
\(Ha(C_1+C_2)\), and the regular budget in
\eqref{eq:ex-lattice-regular-small-budget} pays their sum. The chart
construction has enough room that the support is nonwrapping and every
probe stencil lies in the healthy region; the ordered derivative bounds
through order \(d\) are exactly the ones used by the finite recovery map.
This is the local probe conclusion of the small-support construction, not an
inference from the value bound alone.

For the large-support case, the local field is built by its support
reservation. The reservation includes the factor
\(\exp((\alpha_c+2)q_i)\geq1\), where \(\alpha_c>0\) is the receiving
support exponent and \(q_i\geq0\) the channel's coarse cover number, so
its unweighted supremum cost is no larger than its reserved mass. The
construction bounds that reserved mass using the one-site talent estimate
times the source mass. That estimate is at most \(v/\chi\), hence each
channel has \(S_i\leq(v/\chi)m_0\). The cloud bounds likewise give
\(C_i\leq(v/\chi)m_0\). Applying the local readout bound to the two chart
bounds and the two cloud terms yields
\(OHa(S_1+S_2+C_1+C_2)\leq4OHavm_0/\chi\). The large-support budget pays this
quantity. In both cases the chosen support labels retain the same fine
source and the transition preserves its density, while the one-site
reservation and turn-price inequalities give the remainder \(\rho m_0\).
The local chart statements used in these paragraphs are the exact outputs
of the small-support and large-support constructions; their hypotheses and
producer declarations are recorded in the accompanying evidence file.

In the small-support construction, the image of the source support is first
enlarged by the transport halo. The geometric placement puts this enlarged
set in a parent patch and bounds every point of that patch's support by the
chart radius. The condition
\[
 20(r_d+R_d+1)<n_0
\]
then makes both the transported source support and the parent patch
nonwrapping; \(r_d\) is the chart radius, \(R_d\) the transport reach, and
\(n_0\) the parent side length. The local construction gives the actual
field \(U_i\) and its smooth reference \(V_i\) on one healthy chart. For
each channel weight \(w_i\), and for every ordered list
\((A_1,\ldots,A_\ell)\) of at most \(d\) directions each satisfying the
directional bound of the receiving budget at the origin (entrywise at most
the receiving threshold times the spatial weight), the ordered derivatives
at the identity satisfy
\[
 \bigl|\partial_{A_1}\cdots\partial_{A_\ell}(U_i-V_i)(\mathrm{id})\bigr|
       \leq 2(v/\chi)w_i.
\]
This is a bound on these normalized directions only, not an operator-norm
bound for the derivative.
For every coefficient-recovery probe \(\tau\in\mathcal T\), evaluated at
parameter zero along its test path, one also has
\[
 |\partial^\tau_0 V_i|\leq(v/\chi)w_i.
\]
The probe stencil lies in the healthy chart. The first bound controls the
local variation correction in the reconstruction, and the second is the
probe-price input. The finite recovery map and its probe-price sum therefore
give the \(Wu\chi^{-1}m_0\) contribution; applying the same row decomposition
to the two real channels gives the separate \(Ha(C_1+C_2)\) contribution.
This supplies the stated small-support estimate with the correct one-block
normalization.

In the large-support construction, the transported cover number is larger
than the receiving radius. For each real channel, the support-reservation
argument constructs a regular field with supremum bound \(S_i\) and a
reserved mass
\[
 R_i=e^{(\alpha_c+2)q_i}S_i,
 \qquad R_i\leq(v/\chi)m_0,
\]
where \(\alpha_c>0\) is the receiving support exponent and \(q_i\geq0\)
is the channel's coarse cover number. The second inequality combines the
support-reservation estimate, which bounds \(R_i\) by the tuition of the
selected local record times \(m_0\), with the tuition entry of that
record, which is at most \(v/\chi\). The exponential factor is at least
one, so \(S_i\leq R_i\leq(v/\chi)m_0\). Independently, the cloud estimate
gives \(C_i\leq(v/\chi)m_0\). Thus the two chart bounds and two cloud
charges together cost at most
\[
 OHa(S_1+S_2+C_1+C_2)
       \leq 4OHav\chi^{-1}m_0,
\]
which is paid by the large-support inequality. This is the large-support
estimate used in the proposition; it keeps the chart and cloud contributions
separate because they come from different source bounds.
